For the upper bound, fix an arm \(a\) and estimate \(\mu_a\) from its standalone i.i.d. sample by a boundary-corrected tensor local-polynomial estimator of degree \(\lfloor s\rfloor\), bandwidth \(h\), and projection onto \([-M,M]\). The bounded-density assumption makes every population polynomial Gram matrix uniformly nonsingular, while the Gaussian outcome and variance-range assumptions give a uniform conditional second-moment bound. Taylor expansion of an \(s\)-Hölder function on cubes of side \(h\), followed by the conditional variance calculation, gives
\[
E_P\|\widehat\mu_a-\mu_a\|_{L_2(P_X)}^2
\leq C\{h^{2s}+(nh^d)^{-1}\}.
\tag{40}
\]
Projection cannot increase the error because \(|\mu_a|\leq M\). Taking \(h\asymp n^{-1/(2s+d)}\) and setting \(\widehat g_a=\omega_a\widehat\mu_a\), the known-weight bound yields
\[
E_P\|\widehat g_a-\omega_a\mu_a\|_{L_2(P_X)}^2
\leq W^2 E_P\|\widehat\mu_a-\mu_a\|_{L_2(P_X)}^2
\leq C n^{-2s/(2s+d)}.
\tag{41}
\]
The maximum over the fixed number \(K\) of arms preserves the constant.

For the lower bound, use the active arm \(a_0\), choose the uniform covariate density, and keep every nuisance other than \(\mu_{a_0}\) fixed at an admissible constant. Nonemptiness of the bounded-density class on the unit-volume cube implies \(c_X\leq1\leq C_X\), so this density is admissible. Place disjoint \(s\)-Hölder bumps of width \(h\) and height \(\rho h^s\) on a regular grid inside \(B_0\), with \(\rho>0\) fixed small enough to obey both the Hölder radius and the sup-norm bound. There are \(N_h\asymp h^{-d}\) bumps. For adjacent sign vertices the one-observation Gaussian location divergence is bounded by a constant times \(h^{2s}\) on a region of volume \(h^d\), so the \(n\)-sample Kullback--Leibler divergence is at most \(Cnh^{2s+d}\). Choose \(h\asymp n^{-1/(2s+d)}\) and reduce \(\rho\), if necessary, so every adjacent testing affinity is bounded below by a fixed positive constant. On \(B_0\), \(|\omega_{a_0}|\geq w_{\min}\); hence the weighted squared separation for one changed bit is at least \(c w_{\min}^2h^{2s+d}\). The bit-testing inequality over the \(N_h\) signs gives
\[
\mathsf Q_n^{\mu}(s,d;\omega)
\geq cN_hh^{2s+d}
=c h^{2s}
=c n^{-2s/(2s+d)}.
\tag{42}
\]
No step uses a strict variance interval: one may fix the Gaussian variance at any value allowed by the variance-range assumption. Equations (41)--(42) prove the lemma.